% Generated by scripts/locobench_paper_examples.py from f055.
\paragraph{Question.} How did geologists resolve the dispute over whether the end-Cretaceous mass extinction was triggered by the Chicxulub impact or by Deccan Traps volcanism? (Geology, Natural Sciences.)

\paragraph{Claims.} The plan's claims in assertion order, with their factuality labels; claims joined by a ladder's slots are boxed.

\begin{center}\scriptsize
\setlength{\tabcolsep}{2.5pt}
\begin{tabular}{@{}rp{0.76\linewidth}l@{}}
\toprule
& Claim & Label \\
\midrule
$a_{1}$ & The Alvarez team reported an anomalous enrichment of iridium in the Cretaceous-Paleogene boundary clay at Gubbio, Italy. & true \\
$a_{2}$ & Iridium is scarce in Earth's crust and comparatively abundant in meteoritic material. & true \\
$a_{3}$ & The iridium anomaly has been documented at more than 100 Cretaceous-Paleogene boundary sections worldwide. & true \\
$a_{4}$ & The Alvarez iridium results were published in the journal Nature in 1980. & \textbf{false} \\
$a_{5}$ & Shocked quartz grains are present in the Cretaceous-Paleogene boundary layer. & true \\
$a_{6}$ & Glen Penfield and Antonio Camargo identified a buried circular structure beneath the Yucatán Peninsula during petroleum surveys for Pemex. & true \\
$a_{7}$ & Alan Hildebrand and colleagues published a 1991 paper linking the Chicxulub structure to the Cretaceous-Paleogene boundary. & true \\
$a_{8}$ & The Deccan Traps form a flood basalt province covering much of west-central India. & true \\
$a_{9}$ & The shocked quartz at the boundary layer formed during explosive eruptions of the Deccan Traps. & \textbf{false} \\
$a_{10}$ & Gerta Keller argued that Deccan Traps volcanism was the primary driver of the end-Cretaceous extinction. & true \\
$a_{11}$ & The Chicxulub impact and the Cretaceous-Paleogene boundary are coeval within radiometric dating uncertainty. & true \\
$a_{12}$ & The Chicxulub impact predates the Cretaceous-Paleogene boundary by roughly 300,000 years. & true \\
\slotclaim{$a_{13}$} & IODP-ICDP Expedition 364 drilled into the Chicxulub peak ring in 2016. & true \\
\slotclaim{$a_{14}$} & The Expedition 364 drill core recovered shocked granite and suevite from the Chicxulub peak ring. & true \\
\bottomrule
\end{tabular}
\end{center}

\paragraph{Template.} The answer as written once, with its slot markers.

\begin{quote}\small
The case was assembled measurement by measurement over four decades. The Alvarez team reported an anomalous enrichment of iridium in the Cretaceous–Paleogene boundary clay at Gubbio, Italy. Iridium is scarce in Earth's crust and comparatively abundant in meteoritic material. The Gubbio measurements came first; only afterwards was the same anomaly documented at more than 100 Cretaceous–Paleogene boundary sections worldwide. The Alvarez iridium results were published in the journal Nature in 1980 — that is, the enrichment measured by that team in the Italian boundary clay appeared as a 1980 Nature paper.

Much of the dispute turned on what the boundary layer itself contains, and on how the minerals in it could have been made. Shocked quartz grains are present in the Cretaceous–Paleogene boundary layer, and the petrography of those grains was examined at length by partisans of both the impact and the volcanic hypotheses.

Glen Penfield and Antonio Camargo identified a buried circular structure beneath the Yucatán Peninsula during petroleum surveys for Pemex. More than a decade after that identification, Alan Hildebrand and colleagues published a 1991 paper linking the Chicxulub structure to the Cretaceous–Paleogene boundary. The Deccan Traps form a flood basalt province covering much of west-central India. The shocked quartz at the boundary layer formed during explosive eruptions of the Deccan Traps, a derivation that cannot be squared with the standing of the Deccan Traps as a flood basalt province covering much of west-central India. The eruptive origin of those grains is the evidence on which Gerta Keller argued that Deccan Traps volcanism was the primary driver of the end-Cretaceous extinction. The correlation established in the 1991 Hildebrand paper indicates, for its part, that the Chicxulub impact and the Cretaceous–Paleogene boundary are coeval within radiometric dating uncertainty. Either that indistinguishability of ages holds, or the Chicxulub impact predates the Cretaceous–Paleogene boundary by roughly 300,000 years: the crater is either synchronous with the boundary or some 300,000 years older than it, and it cannot be both.

\slotmarker{S1} Deep drilling into the crater structure thus supplied the physical material against which the competing accounts of the extinction could be tested, and it is on such direct sampling, rather than on argument from correlation alone, that the judgement of the geological community finally rested.
\end{quote}

\paragraph{Fillers.} Each slot's fillers, as admitted by the committee.

\begin{center}\scriptsize
\setlength{\tabcolsep}{2.5pt}
\begin{tabular}{@{}lp{0.80\linewidth}@{}}
\toprule
Filler & Text \\
\midrule
\multicolumn{2}{@{}l}{\slotmarker{S1} control slot, $a_{13}$--$a_{14}$} \\
\quad wording 1 & In 2016 IODP–ICDP Expedition 364 drilled into the Chicxulub peak ring, and as a result the expedition's core recovered shocked granite and suevite from that peak ring. \\
\quad wording 2 & The recovery of shocked granite and suevite from the Chicxulub peak ring came about because IODP–ICDP Expedition 364 drilled into that peak ring in 2016. \\
\quad wording 3 & IODP–ICDP Expedition 364 sank a borehole into the Chicxulub peak ring in 2016, an operation that produced a core of shocked granite and suevite from the peak ring. \\
\bottomrule
\end{tabular}
\end{center}

\paragraph{Rungs.} The five responses of the family, which filler each carries in the edited slots, and gold $\LCS$ at the verdict and the asserted world.

\begin{center}\scriptsize
\setlength{\tabcolsep}{2.5pt}
\begin{tabular}{@{}l>{\raggedright\arraybackslash}p{0.27\linewidth}>{\raggedright\arraybackslash}p{0.30\linewidth}cc@{}}
\toprule
& Ladder, rung & Fillers & \multicolumn{2}{c}{$\LCS$ at world} \\
& & & verdict & asserted \\
\midrule
t0 & order $\rho_0$ & paragraphs 4-3-2-1 & 0.115 & 0.115 \\
t1 & order $\rho_1$, control $\rho_0$ & paragraphs as written & 0.115 & 0.115 \\
t2 & order $\rho_2$ & paragraphs 2-3-4-1 & 0.115 & 0.115 \\
t3 & control $\rho_1$ & S1: wording 2 & 0.115 & 0.115 \\
t4 & control $\rho_2$ & S1: wording 3 & 0.115 & 0.115 \\
\bottomrule
\end{tabular}
\end{center}

\paragraph{Response t1 in full} (order $\rho_1$; edited slots highlighted, 399 words).

\begin{quote}\small
The case was assembled measurement by measurement over four decades. The Alvarez team reported an anomalous enrichment of iridium in the Cretaceous–Paleogene boundary clay at Gubbio, Italy. Iridium is scarce in Earth's crust and comparatively abundant in meteoritic material. The Gubbio measurements came first; only afterwards was the same anomaly documented at more than 100 Cretaceous–Paleogene boundary sections worldwide. The Alvarez iridium results were published in the journal Nature in 1980 — that is, the enrichment measured by that team in the Italian boundary clay appeared as a 1980 Nature paper.

Much of the dispute turned on what the boundary layer itself contains, and on how the minerals in it could have been made. Shocked quartz grains are present in the Cretaceous–Paleogene boundary layer, and the petrography of those grains was examined at length by partisans of both the impact and the volcanic hypotheses.

Glen Penfield and Antonio Camargo identified a buried circular structure beneath the Yucatán Peninsula during petroleum surveys for Pemex. More than a decade after that identification, Alan Hildebrand and colleagues published a 1991 paper linking the Chicxulub structure to the Cretaceous–Paleogene boundary. The Deccan Traps form a flood basalt province covering much of west-central India. The shocked quartz at the boundary layer formed during explosive eruptions of the Deccan Traps, a derivation that cannot be squared with the standing of the Deccan Traps as a flood basalt province covering much of west-central India. The eruptive origin of those grains is the evidence on which Gerta Keller argued that Deccan Traps volcanism was the primary driver of the end-Cretaceous extinction. The correlation established in the 1991 Hildebrand paper indicates, for its part, that the Chicxulub impact and the Cretaceous–Paleogene boundary are coeval within radiometric dating uncertainty. Either that indistinguishability of ages holds, or the Chicxulub impact predates the Cretaceous–Paleogene boundary by roughly 300,000 years: the crater is either synchronous with the boundary or some 300,000 years older than it, and it cannot be both.

\slotfill{S1}{In 2016 IODP–ICDP Expedition 364 drilled into the Chicxulub peak ring, and as a result the expedition's core recovered shocked granite and suevite from that peak ring.} Deep drilling into the crater structure thus supplied the physical material against which the competing accounts of the extinction could be tested, and it is on such direct sampling, rather than on argument from correlation alone, that the judgement of the geological community finally rested.
\end{quote}

\paragraph{Gold relations of response t1.} Sense, compiled coupling, strength band, and the relation's sign at the verdict (v.) and asserted (a.) world ($+$ support, $-$ conflict, $\cdot$ no factor); a slot id marks a relation drawn by a filler.

\begin{center}\scriptsize
\setlength{\tabcolsep}{2.2pt}
\begin{tabular}{@{}lllllcc@{}}
\toprule
Pair & & Sense & Coupling & Band & v. & a. \\
\midrule
$a_{1}\to a_{3}$ &  & Precedence & none & strong & $\cdot$ & $\cdot$ \\
$a_{4}\to a_{1}$ &  & Restatement & equivalence & moderate & $-$ & $+$ \\
$a_{6}\to a_{7}$ &  & Precedence & none & strong & $\cdot$ & $\cdot$ \\
$a_{7}\to a_{11}$ &  & Evidence & entailment & moderate & $+$ & $+$ \\
$a_{8}\to a_{9}$ &  & Refutation & contradiction & moderate & $+$ & $-$ \\
$a_{9}\to a_{10}$ &  & Evidence & entailment & weak & $+$ & $+$ \\
$a_{11}\to a_{12}$ &  & Alternative & exclusive & strong & $-$ & $-$ \\
$a_{13}\to a_{14}$ & S1 & Cause-Effect & entailment & strong & $+$ & $+$ \\
\bottomrule
\end{tabular}
\end{center}
